\specialsection{Список сокращений и обозначений}

\term{API}{программный интерфейс взаимодействия компонентов системы}
\term{UI}{пользовательский интерфейс веб-приложения}
\term{БД}{база данных системы}
\term{Рабочее пространство}{логическая область системы, объединяющая участников, проекты, задачи, документы, файлы и настройки доступа}
\term{Проект}{рабочая область внутри рабочего пространства, в рамках которой ведутся задачи, документация, заметки и история активности}
\term{Задача}{единица работы, имеющая статус, приоритет, исполнителя, сроки, связи и историю изменений}
\term{Группа задач}{крупный логический блок работ, объединяющий связанные задачи проекта}
\term{Связь задач}{направленная зависимость между задачами, на основе которой определяются блокировки и порядок выполнения работ}
\term{Документ}{страница базы знаний проекта с содержимым в формате Markdown}
\term{Файловый ресурс}{загружаемый файл или изображение, связанное с задачами, документами, заметками или другими сущностями системы}
\term{Заметка}{краткая текстовая запись, связанная с проектным контекстом, задачей, документом или агентной сессией}
\term{Агент}{программный участник системы, выполняющий действия через выделенные учетные данные и интеграционный интерфейс}
\term{Учётные данные агента}{ключ доступа агента с ограничением по области разрешений, сроку действия и состоянию}
\term{Сессия агента}{отдельный запуск или рабочая сессия агента, связанная с журналом действий и созданными сущностями}
\term{Интеграция}{внешнее подключение системы к стороннему сервису, например к платформе GitHub}
\term{MCP}{Model Context Protocol, протокол взаимодействия внешних инструментов и агентов с системой}
\term{Аудит}{журнал значимых действий пользователей, агентов и интеграций в системе}